Long-form text generated by large language models (LMs) has widely been used~\citep{gpt3,ouyang2022training};
nonetheless, evaluating their {\em factual precision}---whether each piece of information conveyed in a generation is factually accurate---remains challenging for two reasons.
First, a generation consists of a large number of pieces of information that are a mixture of true or false,\footnote{
Even a single sentence consists of multiple pieces of information (e.g., 4.4 per sentence in ChatGPT, 40\% of which are a mixture of supported and unsupported information).} making a binary judgment inadequate~\citep{pagnoni-etal-2021-understanding}.
% Prior work adds a partial support label~\citep{manakul2023selfcheckgpt,liu2023evaluating} which does not inherently address the issue.
% \pw{why does it not inherently address the issue? This sentence feels a bit out of place in the footnote; perhaps consolidate this in related work (or move it to the main text and expand on it a bit if you think it's important to address in the intro)?} \kalpesh{+1.. let's keep the first sentence of the footnote and move the prior work part to the first two paragraphs of Section 2? Also, we could cite~\citep{pagnoni-etal-2021-understanding} in support of going beyond binary labels}
%\mohit{it's still unclear to me what the challenge is here. let's say a generation has a bunch of info, some of which is true and some of which is false. where is the challenge? is it in validating each piece of info? can you make it more specific?}
Second, validating every piece of information is time-consuming and costly.

\begin{figure}[t]
\centering \footnotesize
\resizebox{\columnwidth}{!}{
    \includegraphics[trim={8.2cm 8cm 14cm 1cm},clip]{figures/teaser.pdf}%
}\vspace{-.3em}
\caption{
    An overview of \ours, a fraction of \emph{atomic facts} (pieces of information) supported by a given knowledge source.
    % \ours\ considers a portion of atomic facts (pieces of information) supported by a given source of knowledge. 
    % \pw{this sounds like \ours only looks at a portion of the output, rephrase?}\kalpesh{+1, how about "An overview of our metric \ours, which checks what fraction of \emph{atomic facts} (pieces of information) are supported by a given knowledge source"}
    %
    \ours\ allows a more fine-grained evaluation of factual precision, e.g., in the figure, the top model gets a score of 66.7\% and the bottom model gets 10.0\%, whereas prior work would assign 0.0 to both.
    % In this figure, factual precision is \textbf{57.14\%}.
    \ours\ can either be based on human evaluation, or be automated, which allows evaluation of a large set of LMs with no human efforts.
}\label{fig:teaser}
\end{figure}

In this paper, we introduce \textbf{\ours} (\ourslong), a new evaluation of an LM that represents {\em the percentage of atomic facts (pieces of information) supported by a given knowledge source}.
% \pw{if this is going to be the official definition, should we use ``atomic facts'' instead of ``pieces of information'' for consistency?} generated by the LM that are supported by a given knowledge source.''
% \pw{nit but the double quotes makes me think this is a quotation of someone else, instead of our own definition}
%
Computing \ours\ involves
%to \pw{grammar}
(1) breaking a generation into a series of atomic facts---short statements that each contain one piece of information~\citep{nenkova2004evaluating,shapira-etal-2019-crowdsourcing,zhang-bansal-2021-finding,liu2022revisiting}, and (2) assigning a binary label %\mohit{binary label?}
% \pw{precision, in the precision/recall sense, is a property of a set of outputs as opposed to a single output, so ``binary precision label'' might be confusing; though perhaps we're just using precision colloquially?} \kalpesh{How about "binary precision label" --> "binary label"? we have precision later on in the sentence}
to each atomic fact, allowing a fine-grained evaluation of factual precision.
%
%We propose to use this pipeline by prompting the LM to generate \emph{people biographies} \pw{maybe consider cleanly separating \ours (as an evaluation metric that involves decomposition into atomic facts) from the fact that in this paper we compute \ours based on human biography generations, since \ours could be used for other types of generations as well?}\kalpesh{+1, but I think a minor rewording will do the job. "We propose to use this pipeline" --> "We evaluate FactScore on the task of human biography generation...."}, because
We evaluate \ours\ on the task of generating people biographies because
% This definition provides reasonable evaluation when LMs are prompted to generate {\em human biographies} \mohit{is the definition only reasonable for biographies? perhaps it's better to say something like ``We use this pipeline to evaluate the factual precision of LM-generated \emph{human biographies}. We choose human biographies because ...}
%and evaluated against Wikipedia, because 
generations
consist of verifiable statements rather than debatable or subjective ones,
%is factoid rather than vague or debatable
and the scope is broad
%enough \pw{-enough}
(i.e., covering diverse nationalities, professions, and levels of rarity). %, and Wikipedia has high coverage. 
%\pw{does `factoid' have a technical meaning? \url{https://en.wikipedia.org/wiki/Factoid} doesn't seem to match how we're using it here} \kalpesh{interesting! I guess NLP folks use factoid the wrong way haha.. how about "generation is factoid" --> "generations consist of a series of verifiable factual statements rather than debatable sentences."}


%We evaluate the factual precision of the state-of-the-art, commercial LMs---InstructGPT, ChatGPT, and search-incorporated, PerplexityAI,\footnote{\label{fn:ppl}\url{https://www.perplexity.ai/}} referred to as \subjectLM{}s.
%
We perform extensive human annotations to obtain \ourScores\ of three state-of-the-art, commercially available LMs: InstructGPT~\citep{ouyang2022training}, ChatGPT~\citep{chatgpt}, and search-augmented PerplexityAI.\footnote{\label{fn:ppl}\href{https://www.perplexity.ai/}{\texttt{perplexity.ai}}}
% We recruit fact-checking experts to validate atomic facts for 505 model generations, which we confirm to be expensive (76 hours of work and \$2K in total). 
Our results indicate that commercially available LMs are riddled with errors, having \ourScores\ of 42\%, 58\% and 71\%, respectively.
Their \ourScores\ significantly drop as the rarity of the entities increases, e.g., $80\%\rightarrow16\%$ for ChatGPT.

Since human evaluation is costly, we next introduce an automatic evaluation of \ours\ through a model that estimates a \ourScore\ for a given LM.
Our estimator decomposes generations into atomic facts and validates each based on a given knowledge source, leveraging retrieval from the given knowledge source and strong language models.
Our estimator closely approximates \ours\ with an error rate of $<2\%$
% Our ablations show that incorporating the knowledge source through retrieval and breaking a generation into a series of atomic facts are key to the accuracy of \oursauto.\mohit{it's not clear from intro how the decomposition into atomic facts is done}
%
and can be applied to a range of {\em new} LMs at scale with no human effort.
% \pw{This could seem circular: are we arguing that a retrieval-augmented LM can accurately assess factual precision even though if we take generations from a retrieval-augmented LM, those generations are themselves not factually precise? Perhaps worth addressing?} \kalpesh{imo not necessary for introduction.. for addressing later in paper (say Sec 4.1.1) we could cite something like GPTScore where GPT models are themselves used to evaluate generations? } \sewon{sg to me}
Our case study evaluates 6,500 generations from 13 LMs that could have cost \$26K, %, and providing new insights about them.
with various findings: %for instance,
GPT-4~\citep{OpenAI2023GPT4TR} and ChatGPT are far less factual than humans % \pw{far worse/less factual than humans?}
but are much better than public models,
%\pw{in the previous paragraph, didn't we already show that ChatGPT isn't doing too well? is this a different result? also, is ChatGPT considered a public model (since it's publicly accessible via API)?}, \sewon{oh I meant openai models aren't open. reworded to clarify this.}
and there is a large variance between public models, with % e.g., \pw{/with}
Vicuna~\citep{vicuna2023} and Alpaca~\citep{alpaca} being some of the best.
% \kalpesh{maybe one sentence her summarizing main results takeaway from the new LM evaluation? something like "smaller public LMs far lag behind ChatGPT/GPT4 in their FSEstimator scores"}

%Moreover, outputs from \ours\ significantly improve an editing system that provides a revision to the factual errors, showing promising results on automated text generation and revisions.


In summary, our contributions are as follows.
\vspace{-.3em}
\begin{enumerate}[leftmargin=15pt]\itemsep -.2em
    \item We introduce \ours, a new evaluation of factual precision of LMs by breaking their generations into atomic facts and validating each against a given knowledge source. Human evaluation reveals that the state-of-the-art LMs with and without search have low \ourScores. %\ (42.5--71.5\%).
    \item We introduce a model that approximates \ours\ with an error rate of $<2\%$, allowing evaluation of a large set of new LMs without manual human efforts.
    \item We open-sourced \ours\ and the annotated data for public use, available via \texttt{pip install factscore}.
    We suggest future work to extend \ours\ for a broader set of generations (e.g., open-ended generation) and to further improve the estimator.
\end{enumerate}
\pw{Overall, I think the introduction could frame this paper more clearly in the context of prior work: what is the key idea and how is it different from prior work? And what are the results/analysis/insights that this idea enables?}